\section{Application: Wind Farm Layout Optimization}
\label{sec:application}

\subsection{Wake Model}
\label{sec:wake}

We use the Bastankhah Gaussian wake deficit model~\cite{bastankhah2014new}, which describes the velocity deficit behind a wind turbine as a Gaussian profile:
\begin{equation}
\label{eq:deficit}
    \frac{\Delta u}{u_\infty} = C(x, C_T) \cdot \exp\!\left( -\frac{r^2}{2\sigma(x)^2} \right)
\end{equation}
where $\Delta u$ is the velocity deficit at downstream distance $x$ and radial offset $r$, $u_\infty$ is the freestream wind speed, $C_T$ is the thrust coefficient, and $C(x, C_T)$ is the centerline deficit derived from momentum conservation.
The wake width $\sigma$ grows linearly with downstream distance:
\begin{equation}
\label{eq:sigma}
    \sigma(x) = k \, x + \epsilon_0 \, D
\end{equation}
where $k$ is the wake expansion coefficient (fixed at $k = 0.04$ for all experiments), $D$ is the rotor diameter, and $\epsilon_0$ is the near-wake expansion parameter derived from the thrust coefficient.

The wake simulation is implemented in \texttt{pixwake}, a JAX-native wake model library that supports automatic differentiation through the entire power computation.
Wake superposition uses a fixed-point iteration with a custom VJP (vector-Jacobian product) rule, enabling \texttt{jax.grad} to compute exact gradients of AEP with respect to turbine positions.
The LLM's optimizer receives the simulation object and can compute gradients freely, but cannot modify the wake model parameters.

\subsection{Training Farm: Danish Energy Island (DEI)}
\label{sec:dei}

The training farm is based on Danish Energy Island tender area farm~1~\cite{energinet2022dei}.
Key parameters:
\begin{itemize}[nosep]
    \item $N = 50$ turbines, IEA 15\,MW reference turbine~\cite{gaertner2020iea15mw}, $D = 240$\,m, hub height 150\,m.
    \item $s_{\min} = 960$\,m ($4D$).
    \item 5-vertex convex polygon boundary ($\sim$20\,km $\times$ 35\,km).
    \item 24-sector wind rose derived from 10 years of observed wind data (daily averages), with sector-specific mean wind speeds ranging from 7.3 to 11.5\,m/s.
\end{itemize}
The polygon is spacious relative to the turbine count: 50 turbines at 960\,m spacing occupy a small fraction of the available area.
This means that boundary and spacing constraints are relatively easy to satisfy, a property that matters for understanding the constraint-handling failure mode we observe.

\subsection{Held-Out Farm: IEA Wind ROWP}
\label{sec:rowp}

The held-out farm is the IEA Wind Task~37 Reference Offshore Wind Plant (ROWP) irregular layout case~\cite{borrassunyeriea2024rowp}.
Key parameters:
\begin{itemize}[nosep]
    \item $N = 74$ turbines, IEA 10\,MW reference turbine~\cite{bortolotti2019iea10mw}, $D = 198$\,m, hub height 119\,m.
    \item $s_{\min} = 792$\,m ($4D$).
    \item 4-vertex convex polygon boundary ($\sim$19\,km $\times$ 16\,km), a parallelogram-like shape more constrained than the DEI polygon.
    \item 12-sector Weibull wind rose with mean wind speeds from 7.2 to 11.3\,m/s.
\end{itemize}

The ROWP farm differs from the training farm along every axis that an optimizer might exploit: the turbine count is 48\% larger, the rotor diameter is 17.5\% smaller, the minimum spacing is 17.5\% tighter, the polygon has fewer vertices and a more elongated aspect ratio, and the wind rose has half as many sectors with a different frequency distribution.
These differences are not artificial---they reflect real variation across offshore wind planning cases.

\subsection{Constraint Handling in the Baseline Solver}
\label{sec:constraints}

The \texttt{topfarm\_sgd\_solve} optimizer handles constraints via differentiable penalty functions added to the objective.
The augmented loss at iteration $t$ is:
\begin{equation}
\label{eq:augmented}
    \mathcal{L}_t(\mathbf{x}, \mathbf{y}) = f(\mathbf{x}, \mathbf{y}) + \alpha_t \left[ w_b \cdot g_b(\mathbf{x}, \mathbf{y}) + w_s \cdot g_s(\mathbf{x}, \mathbf{y}) \right]
\end{equation}
where $f$ is the negative AEP objective, $g_b$ is the boundary penalty (sum of squared boundary violations), $g_s$ is the spacing penalty ($\sum \max(0, s_{\min}^2 - d_{ij}^2)$ for all pairs), and $w_b$, $w_s$ are user-specified penalty weights.

The penalty multiplier $\alpha_t$ follows an adaptive schedule tied to the learning rate decay:
\begin{equation}
\label{eq:alpha}
    \alpha_t = \alpha_0 \cdot \frac{\eta_0}{\eta_t}
\end{equation}
where $\alpha_0$ is initialized from the gradient magnitude at iteration 0 ($\alpha_0 = \overline{|\nabla f_0|} / \eta_0$), $\eta_0$ is the initial learning rate, and $\eta_t$ decays according to $\eta_t = \eta_{t-1} / (1 + \gamma \cdot t)$.
As the learning rate decreases, the constraint penalty increases, ensuring that the optimizer converges to a feasible point.
This mechanism is critical: it allows large infeasible moves early (when the learning rate is high) while enforcing feasibility as the optimizer settles.

This adaptive ramping is the mechanism that no LLM-generated custom optimizer reproduces.
Custom optimizers that use fixed penalties ($\alpha_t = \text{const}$) or decreasing penalties work on the spacious DEI farm but fail when the ROWP polygon leaves less margin for constraint violation.

\subsection{Models and Compute}
\label{sec:models}

We report primary results using Gemini~2.5~Flash~\cite{gemini2025flash} through the Gemini API with function-calling mode.
The agent architecture (Section~\ref{sec:agent}) is backend-agnostic: we additionally run experiments with Claude (via Claude Code CLI) and Llama~3.1~405B~\cite{dubey2024llama3} (self-hosted via vLLM~\cite{kwon2023vllm}).

\paragraph{Time budget.}
Each run has a 5-hour wall-clock time budget.
Individual optimizer evaluations time out at 60 seconds during the iterative search phase and 300 seconds for the final evaluation.
The 60-second timeout prevents multi-start strategies with large start counts from consuming the time budget during exploration.

\paragraph{Hot-start.}
All runs begin from the same seed optimizer: a single-start \texttt{topfarm\_sgd\_solve} with default hyperparameters (learning rate~$=50$, penalty weights~$=1.0$).
This seed scores below the 500-start baseline, providing a working template without giving away the strategies that improve over it.

\paragraph{Reproducibility.}
Wake simulations are deterministic (JAX with 64-bit floating point).
LLM sampling introduces stochasticity; we plan multiple seeds per model to report confidence intervals.
All optimizer scripts are versioned, and the full attempt log with paired train/held-out scores is saved for post-hoc analysis.
